\scanpage{290}
\mspage{240}

ersion ouverte de Dwyer), et $p_f\in W$.

7. \uline{Retour sur l'axiomatique de la fibre $W$-homotopique}.

Je reviens sur \unsure{les réflexions|la réflexion} du n\textsuperscript{o} 4. Finalement, pour
\[
 (M,W\subset Fl\,M)
\]
\unsure{donnés}, avec $W$ faiblement saturé, il apparaît qu'une excellente théorie des fibres $W$-homotopiques, satisfaisant aux conditions \textcircled{a} \textcircled{b} \textcircled{c} du n\textsuperscript{o} 4 (page 215-218), existe \unsure{déjà|dès|sitôt} que l'axiome \deleted{\unsure{anodin|ci-dessous}} d'apparence anodine suivant est satisfait (en plus de l'existence \inserted{de l'objet final et} des produits fibrés)

% slanted
\margin{et encore ! Il me semble qu'il suffit de supposer que les $X\to Y$ qui sont $W$-\unsure{fibrés} sont ``\unsure{quarrables}''.}

\uline{Factorisation gauche} : Toute flèche \inserted{$f:\unsure{\alpha|X}\unsure{-|\to}Y$} de $M$ se factorise en
\[
 \begin{aligned}
 X&\xrightarrow{i_f}\overline X\xrightarrow{p_f}Y,\\
 i_f&\in W,\quad p_f\in Fib_W.
 \end{aligned}
\]
% slanted
\margin{\uline{NB} L'axiome de \deleted{\unsure{suffisant|factorisation}} factorisation \unsure{à gauche|gauche} n'est suffisant pour l'existence d'une bonne théorie des fibres homotopiques \inserted{(\unsure{semble-t-il|malentendu})}, que dans le cas où tout $X$ de $M$ est $W$-fibrant, i.e. est dans $\underline{Fib}_W\,e$.}

% slanted
\margin{Sous ces conditions $M$ muni de $(W,F_W)$ est une \unsure{catégorie|catég.} à \unsure{fibrations|cofibrations} de K.S. \deleted{\unsure{Br.}} Brown -- \unsure{Au dessous|Au dessus} \unsure{de|cf.} ``tout est OK'' mais il faudrait vérifier a) b) c) ci-dessus (p. 215, 218) je ne suis pas sûr de l'avoir fait}

L'axiome dual donne lieu à une bonne théorie des cofibres $W$-homotopiques :

\uline{Factorisation droite} : Toute flèche $f:X-Y$ dans $M$ se factorise en
\[
 X\xrightarrow{j_f}\widetilde Y\xrightarrow{q_f}Y,
\]
avec $q_f\in W$, $j_f\in Cofib_W$.

% archive sheet number: 290

